% BHU PYQ -- Manually transcribed & error-corrected
% Paper: BPT-101 : Mechanics and Relativity
% Exam: B.Sc. (Hons) Semester I Examination, 2022-23

\documentclass[11pt,a4paper]{article}
\usepackage[a4paper,top=2.5cm,bottom=2.5cm,left=2.8cm,right=2.8cm,headheight=14pt]{geometry}
\usepackage{amsmath,amssymb}
\usepackage[T1]{fontenc}
\usepackage[utf8]{inputenc}
\usepackage{lmodern}
\usepackage{microtype}
\usepackage{fancyhdr}
\usepackage{enumitem}
\usepackage{hyperref}
\usepackage{lastpage}
\usepackage{parskip}

\hypersetup{colorlinks=true,urlcolor=black,linkcolor=black,
  pdftitle={BPT-101: Mechanics and Relativity -- BHU PYQ},pdfauthor={Banaras Hindu University}}

\pagestyle{fancy}\fancyhf{}
\renewcommand{\headrulewidth}{0.4pt}\renewcommand{\footrulewidth}{0.4pt}
\fancyhead[L]{\small   \textit{Banaras Hindu University}}
\fancyhead[R]{\small   \textit{BPT-101: Mechanics and Relativity}}
\fancyfoot[C]{\small Page  \thepage\ of \pageref{LastPage}}


\newlist{parts}{enumerate}{2}
\setlist[parts,1]{label=(\alph*),leftmargin=*,itemsep=5pt,topsep=3pt}
\setlist[parts,2]{label=(\roman*),leftmargin=*,itemsep=3pt,topsep=2pt}

\newcommand{\pts}[1]{\hfill   \textit{[#1]}}


\begin{document}


\begin{center}
  {\large\bfseries BANARAS HINDU UNIVERSITY}\\[2pt]
  {\normalsize B.Sc. (Hons.) --- Semester I Examination, 2022-23}\\[6pt]
  \rule{0.8\linewidth}{0.5pt}\\[6pt]
  {\large\bfseries Paper No. BPT-101: Mechanics and Relativity}\\[4pt]
  {\normalsize Time: 3 Hours \hfill Maximum Marks: 70}
\end{center}
\rule{\linewidth}{0.4pt}
\smallskip



\noindent   \textit{Instructions: Answer   \textbf{five} questions including Question~1 which is compulsory. Symbols have their usual meanings.}
\bigskip




\noindent   \textbf{Question 1.}   \textit{(Compulsory --- 3.5 marks each for (c)-(f), 3 marks each for (a)-(b). Answer any four parts.)}

\begin{parts}
  \item Find the amount of work done to increase the speed of an electron from $0.6c$ to $0.8c$. The rest mass of an electron is $0.5\, \text{MeV}$. \pts{3}
  \item The volume of a solid does not vary with pressure. Find Poisson's ratio for the solid. \pts{3}
  \item Explain why a cantilever of uniform cross-section is more likely to break near the fixed end. \pts{3.5}
  \item Calculate the fictitious force and observed force acting on a body of mass $6\, \text{kg}$ kept in frames of reference accelerated vertically upward and downward with acceleration $490\, \text{cm/s}^2$ and $9.8\, \text{m/s}^2$ respectively (take $g = 9.8\, \text{m/s}^2$). \pts{3.5}
  \item Explain the phenomenon of twin paradox. Is aging reciprocal? \pts{3.5}
  \item How much faster than its present speed should the Earth rotate so that bodies lying on the equator may fly off into space? \pts{3.5}
\end{parts}

\medskip\rule{\linewidth}{0.2pt}

\medskip


\noindent   \textbf{Question 2.}

\begin{parts}
  \item Find the twisting couple per unit twist in the case of a hollow cylinder whose inner radius is $R_1$ and outer radius is $R_2$. \pts{5}
  \item Derive the relation connecting Young's modulus ($Y$), Bulk modulus ($K$) and Poisson's ratio ($\sigma$) of a material. \pts{5}
  \item Calculate the mass of water flowing in 2 minutes through a tube $0.2\, \text{cm}$ in inner diameter and $40\, \text{cm}$ long, if there is a constant pressure head of $40\, \text{cm}$ of water. The coefficient of viscosity of water is $0.0085$ CGS units. \pts{4}
\end{parts}

\medskip\rule{\linewidth}{0.2pt}

\medskip


\noindent   \textbf{Question 3.}

\begin{parts}
  \item Show that for a compound pendulum the centers of suspension and oscillation can be interchanged. Explain how a compound pendulum is superior to a simple pendulum. \pts{6}
  \item Derive the equation of motion of a moving coil galvanometer and discuss the conditions under which the motion is (i) deadbeat, (ii) critical and (iii) ballistic. \pts{8}
\end{parts}
\hfill   \textbf{OR}

\begin{parts}
  \item Two particles of masses $m_1$ and $m_2$ collide and after collision move at angles $ heta_1$ and $ heta_2$ respectively. Show that in the centre of mass frame, the velocities of the particles do not change in magnitude after collision. \pts{8}
  \item What are the differences between elastic and inelastic collisions? Write down one example for both types of collisions. \pts{6}
\end{parts}

\medskip\rule{\linewidth}{0.2pt}

\medskip


\noindent   \textbf{Question 4.}

\begin{parts}
  \item In the experiment of meson decay, how are the observations in the Lab frame and the particle frame explained? Take the half-life of mesons = $2.2\,\mu \text{s}$ and velocity = $0.99c$. \pts{7}
  \item Let the velocity of a particle be $v$. Discuss the energy distribution in three conditions: (i) $v \ll c$, (ii) $v \sim c$ and (iii) $v \approx c$. \pts{7}
\end{parts}
\hfill   \textbf{OR}

\begin{parts}
  \item Consider the relativistic form of Newton's second law of motion. Show that when the force $\vec{F}$ is parallel to the velocity $\vec{v}$:
  \[
  F = m_0 \left(1 - \frac{v^2}{c^2}\right)^{-3/2} \frac{dv}{dt}
  \]
  where $m_0$ is the rest mass of the object and $v$ is its speed. \pts{10}
  \item Two spaceships approach each other, moving with the same speed as measured by a stationary observer on the Earth. Their relative speed is $0.70c$. Determine the velocity of each spaceship as measured by a stationary observer on the Earth. \pts{4}
\end{parts}

\medskip\rule{\linewidth}{0.2pt}

\medskip


\noindent   \textbf{Question 5.}

\begin{parts}
  \item What is the Coriolis force? Under what conditions does it come into play? \pts{4}
  \item Discuss in general terms the effect of the Coriolis force produced as a result of the Earth's rotation. \pts{6}
  \item Calculate the values of the centrifugal and the Coriolis forces on a mass of $20\, \text{g}$ placed at a distance of $10\, \text{cm}$ from the axis of a rotating frame of reference, if the angular speed of rotation of the frame is $10\, \text{rad/s}$. \pts{4}
\end{parts}

\vfill
\rule{\linewidth}{0.4pt}

\begin{center}\small
  \href{https://bhu-exam.vercel.app/}{Click here to visit our website}\\[4pt]
  \href{https://me-aryan.vercel.app/}{Click here to visit our contributor}\\[4pt]
  \href{https://quashan-me.vercel.app/}{Click here to visit our contributor}
\end{center}

\end{document}
